
\section{Discussion and Limitations}
\label{sec:discussion}

\subsection{Interpretation of Results}

Table~\ref{tab:writing} illustrates the proposed paired
analysis using four fictional participants. The synthetic
mean completion time was 23.38 minutes without AI and
19.50 minutes with AI, corresponding to a mean paired
time saving of 3.875 minutes. However, the 95\% bootstrap
confidence interval was [0.00, 6.25] minutes and included
zero. More importantly, these observations were
synthetically generated. They validate the statistical
workflow but cannot establish a real productivity effect.

The demonstration also produced higher average quality
scores and required-fact coverage under the AI-assisted
condition. Nevertheless, the proportion of unsupported
claims was higher. This illustrates why completion time,
quality, and factual correctness must be assessed together.

Previous research has examined the effects of generative
AI on professional tasks~\cite{noyExperimentalEvidenceProductivity2023,
dellacquaNavigatingJaggedTechnological2026}.
Our synthetic results should not be directly compared
with empirical effect sizes from those studies.

Table~\ref{tab:software} summarizes a separate benchmark
of 100 Python analysis executions. These measurements
evaluate software performance and logging reliability,
not human productivity or AI inference speed.

\subsection{Limitations and Threats to Validity}

\textit{Internal Validity:}
No real participants were recruited, and the four
participant pairs are fictional. Consequently, the
demonstration cannot support causal conclusions.
A future experiment must control task-order effects,
participant learning, fatigue, and reviewer bias.

\textit{External Validity:}
The proposed study focuses on two matched professional
writing tasks. Their equivalence has not yet been
established through pilot testing. Findings from a future
study may not generalize to other professions or tasks.
Software execution times may also vary across hardware
and operating systems.

\textit{Construct Validity:}
Completion time alone does not capture the full value
of professional work. Quality assessment requires
consistent scoring criteria and trained reviewers.
The proposed five-point non-inferiority margin remains
provisional and requires statistical justification.
Furthermore, Python execution time is not a substitute
for human task completion time.

These limitations define the boundaries of the current
work. The primary contribution is a reproducible
research protocol and analysis pipeline, rather than
an empirical demonstration of AI effectiveness.
